\documentclass[11pt,a4paper]{article}

\usepackage[utf8]{inputenc}
\usepackage[T1]{fontenc}
\usepackage{amsmath,amssymb,amsfonts,amsthm}
\usepackage{mathtools}
\usepackage{geometry}
\usepackage{booktabs}
\usepackage{hyperref}
\usepackage{enumitem}
\usepackage{array}

\geometry{margin=1in}
\hypersetup{
    colorlinks=true,
    linkcolor=blue,
    citecolor=blue,
    urlcolor=blue,
    pdftitle={The Fool's Theorem},
    pdfauthor={Shaurya Sanyal}
}

\theoremstyle{plain}
\newtheorem{theorem}{Theorem}
\newtheorem{proposition}{Proposition}
\newtheorem{lemma}{Lemma}
\newtheorem{corollary}{Corollary}

\theoremstyle{definition}
\newtheorem{definition}{Definition}
\newtheorem{assumption}{Assumption}
\newtheorem{observation}{Observation}

\theoremstyle{remark}
\newtheorem{remark}{Remark}

\newcommand{\Sinv}{\mathcal{S}_{\mathrm{inv}}}
\newcommand{\Snoise}{\mathcal{S}_{\mathrm{noise}}}
\newcommand{\Swork}{\mathcal{S}_{\mathrm{work}}}
\newcommand{\Snon}{\mathcal{S}_{\mathrm{noninv}}}
\newcommand{\Ainv}{A_{i,\mathrm{inv}}}

\title{%
  \textbf{The Fool's Theorem}\\[0.4em]
  \large Error Gravitation in Long-Horizon Agents:\\
  Attention Dilution, the Dilution Threshold, and Bounded Context as a Remedy
}

\author{%
  \textbf{Shaurya Sanyal}\\
  \textit{Team Auralis}\\
  \texttt{shauryasanyal@gmail.com}
}

\date{October 2026}

\begin{document}

\maketitle

% ---------------------------------------------------------------------------
\begin{abstract}
\noindent
We give a complete account of a failure mode in long-horizon tool-using agents
that we call \emph{Error Gravitation}, and that we refer to informally as the
\emph{Fool's Theorem}. The name is a label for the phenomenon, not a claim of
mathematical novelty, and \S\ref{sec:name} states that distinction precisely.

The account has three parts. First, we formalize the mechanism. Given an
attention query whose keys split into a fixed invariant set of size $L_u$ and a
growing distractor set of size $N$, the invariant share satisfies an exact
identity
\[
  \Ainv = \sigma\!\left(\Delta + \ln \frac{L_u}{N}\right),
\]
where $\Delta$ is the mean logit gap separating the two sets and $\sigma$ is the
logistic function. Two consequences follow immediately and are sharp rather than
worst-case: $\Ainv > \tfrac12$ if and only if $N < L_u e^{\Delta}$, and bounding
$N \le C$ yields the closed-form guarantee
$\Ainv \ge \bigl(1 + (C/L_u)e^{-\Delta}\bigr)^{-1}$, which requires no
norm-boundedness assumption and is non-vacuous at realistic scale. The
logarithm in the threshold is the entire story: a model must supply a logit gap
growing as $\ln N$ merely to hold the invariant share constant.

Second, we measure this on a real transformer. Using Qwen2.5-0.5B in eager mode
and reading the true final-query attention distribution, we find that the
effective gap follows
\[
  \hat{\Delta}(N) \;\approx\; 0.945\,\ln N - 4.195
  \qquad (R^2 = 0.9996),
\]
a slope indistinguishable from~$1$. The model therefore compensates for context
growth at almost exactly the rate the theory says is necessary, and the observed
dilution is far weaker than the uniform-score analysis predicts: the invariant
share falls only from $0.427$ to $0.388$ across a $32\times$ increase in
distractor volume. We regard this as the single most important finding in the
paper, and it tempers our own remedy: on this model, with this distractor
distribution, the invariant share does not collapse. What volume control buys is
not rescue from collapse but a bound that does not depend on the model happening
to compensate.

Third, we describe and prove the systems response. \emph{Generational Context}
tiers active context by invariance and places a hard ceiling on the mutable
tiers; \emph{Attention-Weighted Eviction} maintains that ceiling as an enforced
post-condition rather than an assumption. Under it, the closed-form guarantee
becomes horizon-independent. We prove the resulting \emph{Generational Retention
Guarantee}, and we state a set of falsifiable predictions that would refute the
theory, together with the experiments that would test them.

The paper is deliberately two-sided. It argues that volume is the variable a
system designer can control, and it documents the evidence that the model was
already partly controlling it on its own.
\end{abstract}

% ---------------------------------------------------------------------------
\section{On the name, and what is being claimed}
\label{sec:name}

The phrase \emph{Fool's Theorem}, as used in this project, names a
\emph{phenomenon}: the tendency of a long-horizon agent to abandon its stated
objective and settle into a loop of local, symptomatic repairs. It is not a
theorem in the mathematical sense, and nothing in this document should be read
as claiming otherwise.

This distinction is not pedantry. The failure mode is real and worth naming.
But the mathematics that explains it is elementary, and the two results this
paper actually proves are named separately:

\begin{itemize}[leftmargin=*]
  \item the \textbf{Softmax Inundation Bound} (\S\ref{sec:bound}), a worst-case
        statement derived from bounded projection norms; and
  \item the \textbf{Dilution Threshold} (\S\ref{sec:threshold}), an exact
        identity relating the invariant attention share to the mean logit gap;
        and
  \item the \textbf{Generational Retention Guarantee} (\S\ref{sec:guarantee}),
        a horizon-independent lower bound under a maintained volume ceiling.
\end{itemize}

The third is the result we would defend under review. The second we consider the
useful one in practice. The first is included because it is the statement most
likely to be quoted, and it is the weakest of the three, so we would rather
present it with its limitations attached than have it circulate without them.

\begin{quote}
\textbf{A note on how to read the empirical section.} The measurements in
\S\ref{sec:empirical} do not confirm the remedy. They partly undercut the
motivation for it, by showing that the model we tested compensates for context
growth far better than a naive analysis predicts. We report this prominently
because a paper that only reported the favourable half of its own evidence would
not be worth the reader's time.
\end{quote}

% ---------------------------------------------------------------------------
\section{The failure mode}
\label{sec:phenomenon}

\subsection{What it looks like}

Long-horizon agents fail in a characteristic way that is distinct from, and
easier to mistake for, ordinary reasoning error. The symptoms are consistent
across domains:

\begin{enumerate}[label=(\alph*), leftmargin=*]
  \item \textbf{Objective drift.} By the middle of a session, the agent is
        optimizing something adjacent to, but not identical with, the stated
        goal. It fixes the error it most recently saw rather than the error that
        caused the run.
  \item \textbf{Symptom loops.} Successive turns apply variations of the same
        local repair. The agent reports progress because the immediate
        observation changed.
  \item \textbf{Recency dominance.} The most recent tool output determines the
        next action, regardless of whether it is diagnostically relevant. A
        misleading error message encountered late in a run outranks a decisive
        one encountered early.
  \item \textbf{Loss of constraint.} Constraints stated at the start
        (``do not restart the service'', ``do not modify the schema'') stop
        constraining behaviour, not because the agent decided to violate them
        but because they have stopped competing for attention.
\end{enumerate}

None of these require the model to be incapable. An agent can be entirely
competent at local reasoning and still exhibit all four, because the failure is
in \emph{what the model is looking at}, not in how well it thinks about what it
looks at.

\subsection{A worked trace}

The following is a condensed illustration, using a production-incident shape
that recurs in practice. Only the parts that matter are shown.

\begin{quote}
\texttt{\small
turn 1 \; USER: Build is failing after deploy. Find the root cause. Do not\\
\phantom{turn 1 }restart services. Do not modify the database schema.\\
turn 1 \; ASSISTANT: I'll start by checking the deploy logs.\\
turn 2 \; TOOL: ERROR connection refused to redis:6379\\
turn 2 \; ASSISTANT: A Redis connectivity failure. Restarting the cache.\\
turn 3 \; TOOL: ERROR connection refused to redis:6379\\
turn 3 \; ASSISTANT: The restart did not help. Restarting the cache again.\\
turn 4 \; TOOL: WARN migration 0042 not applied; column user.locale missing\\
turn 4 \; ASSISTANT: Redis is still unreachable. Increasing the retry budget.\\
turn 5 \; TOOL: ERROR connection refused to redis:6379\\
turn 5 \; ASSISTANT: The cache is the blocker. Restarting it once more.\\
turn 6 \; TOOL: ERROR schema mismatch on user.locale\\
turn 6 \; ASSISTANT: Restarting the cache to clear the schema error.
}
\end{quote}

Turn~4 contains the root cause and turn~6 repeats it, but neither displaces the
Redis failure from the agent's working set. The agent also violates two explicit
constraints from turn~1. It is not that the model cannot read; it is that the
constraint tokens and the diagnostic tokens have both stopped winning attention
against a growing body of identical error lines.

Note what makes this example instructive: the distractor is not merely noise. It
is a \emph{plausible} error, which is why it holds attention so effectively. The
failure mode is worst when the noise is superficially actionable.

\subsection{Why the mechanism is not the obvious one}

The natural diagnosis is capacity: the context is full, so the model has
forgotten the goal. That diagnosis is usually wrong, or at least incomplete.
Modern windows are large enough that a trace like the one above occupies a tiny
fraction of them, and the instruction is still present in the input. What has
changed is not availability but \emph{allocation}: the instruction is present
and receives a negligible share of attention mass.

This is the distinction the paper formalizes. It is also why the fix is not
``use a bigger window'': doubling the window doubles the room for distractors as
well.

% ---------------------------------------------------------------------------
\section{Formal setting}
\label{sec:prelim}

\subsection{Transcript and length}

Let an agent interact with an environment over discrete turns $t \in
\mathbb{N}$. Define the initial context length
\begin{equation}
  N_t = L_u + \sum_{k<t}\bigl(L_a^{(k)} + L_o^{(k)} + L_r^{(k)}\bigr),
\end{equation}
where $L_u$ is the length of the invariant block (system prompt, standing
objective, hard constraints), and $L_a^{(k)}, L_o^{(k)}, L_r^{(k)}$ are the
token counts of the action, raw observation, and reasoning trace at turn $k$.

\subsection{Attention}

For query position $i$ and key position $j$ within the same head and layer,
\begin{equation}
  A_{i,j}
  = \frac{\exp(s_{i,j})}{\sum_{m=1}^{N_t}\exp(s_{i,m})},
  \qquad
  s_{i,j} = \frac{\mathbf{q}_i^{\top}\mathbf{k}_j}{\sqrt{d_k}},
  \label{eq:attn}
\end{equation}
with $\mathbf{q}_i = \mathbf{X}_i W_Q$ and $\mathbf{k}_j = \mathbf{X}_j W_K$.
We write $\Ainv = \sum_{j \in \Sinv} A_{i,j}$ for the aggregate mass on the
invariant block.

\subsection{Partition}

Fix $t$ and partition the index set into
\begin{align}
  \Sinv &: \text{invariant tokens}, & |\Sinv| &= L_u, \\
  \Snoise &: \text{raw observations}, & |\Snoise| &= \textstyle\sum_k L_o^{(k)}, \\
  \Swork &: \text{actions and reasoning}, & |\Swork| &= \textstyle\sum_k \bigl(L_a^{(k)}+L_r^{(k)}\bigr),
\end{align}
and write $N = |\Snoise|$ for the distractor count, treating $\Swork$ as part of
the non-invariant pool where stated.

\subsection{What is measured, and what is assumed}

It is worth separating these two cleanly, because the empirical section
measures one quantity and the theorems assume properties of another.

\emph{Measured:} the attention distribution $A_{i,\cdot}$ for the final query
position, read directly from the model in eager mode. No approximation, no
surrogate embedding.

\emph{Assumed (only where stated):} bounds on the projections, used solely for
the worst-case result of \S\ref{sec:bound}. The threshold result of
\S\ref{sec:threshold} requires no such assumption, which is one reason we prefer
it.

% ---------------------------------------------------------------------------
\section{The Dilution Threshold}
\label{sec:threshold}

This section contains the paper's most useful piece of mathematics, and it is
short. The point is to replace a worst-case inequality over unknown norms with
an identity over a single measurable quantity.

\begin{proposition}[Two-group softmax identity]
\label{prop:identity}
Fix a query $i$. Suppose the invariant keys share a common score
$s_{\mathrm{inv}}$ and the non-invariant keys share a common score
$s_{\mathrm{non}}$, with sizes $L_u$ and $N$ respectively. Define the
\emph{effective logit gap}
\begin{equation}
  \Delta = s_{\mathrm{inv}} - s_{\mathrm{non}} .
\end{equation}
Then
\begin{equation}
  \Ainv
  = \frac{L_u e^{s_{\mathrm{inv}}}}{L_u e^{s_{\mathrm{inv}}} + N e^{s_{\mathrm{non}}}}
  = \frac{1}{1 + \dfrac{N}{L_u}e^{-\Delta}}
  = \sigma\!\left(\Delta + \ln\frac{L_u}{N}\right),
  \label{eq:threshold}
\end{equation}
where $\sigma(z) = (1+e^{-z})^{-1}$.
\end{proposition}

\begin{proof}
Factor $e^{s_{\mathrm{inv}}}$ from numerator and denominator and divide through
by $L_u$. The logistic form follows from
$1/(1+e^{-z}) = \sigma(z)$ with
$z = \Delta + \ln(L_u/N)$, since
$\dfrac{N}{L_u}e^{-\Delta} = e^{-\Delta - \ln(L_u/N)}$.
\end{proof}

\begin{corollary}[Half-mass threshold]
\label{cor:half}
$\Ainv > \tfrac12$ if and only if
\begin{equation}
  \Delta > \ln\frac{N}{L_u}
  \qquad\Longleftrightarrow\qquad
  N < L_u e^{\Delta}.
\end{equation}
\end{corollary}

\begin{corollary}[Fractional guarantee under a cap]
\label{cor:cap}
For any $\varepsilon \in (0,1)$, if $N \le C$ then
\begin{equation}
  \Ainv \;\ge\; \frac{1}{1 + \dfrac{C}{L_u}e^{-\Delta}}
  \;=\; \frac{L_u e^{\Delta}}{L_u e^{\Delta} + C},
  \label{eq:cap}
\end{equation}
with equality when $N = C$.
\end{corollary}

\begin{remark}[Why this is better than a norm bound]
\label{rem:better}
Proposition~\ref{prop:identity} contains no unknown constants. Equation
\eqref{eq:cap} is an equality at $N=C$, so it is sharp for the two-group model
rather than merely an upper or lower bound with slack. It requires no boundedness
assumption on the projections. And every quantity in it is measurable:
$\Delta$ is estimable from a handful of forward passes, as we do in
\S\ref{sec:empirical}. By contrast the norm-based bound of \S\ref{sec:bound}
depends on $\kappa = e^{2B}$, which for realistic projection norms exceeds
$10^{8}$, and which is uninformative at any practical context length.
\end{remark}

\begin{remark}[The logarithm is the whole phenomenon]
\label{rem:log}
Corollary~\ref{cor:half} says that the invariant set holds the majority of
attention exactly when the model's logit gap exceeds $\ln(N/L_u)$. The
requirement is logarithmic in the distractor count. This has two consequences
that pull in opposite directions, and both matter:

\begin{enumerate}[label=(\roman*), leftmargin=*]
  \item \textbf{Compensation is cheap.} A model need only increase its
        invariant-versus-background gap by about $\ln N$ nats --- for
        $N = 20{,}000$ and $L_u = 62$, roughly $5.8$ nats --- to hold the share
        constant. Since logits are unbounded and gaps of this size are well
        within the range trained models routinely produce, we should \emph{expect}
        capable models to compensate, at least partially. \S\ref{sec:empirical}
        finds that the model we tested does exactly this.
  \item \textbf{Dilution is not avoidable by scale.} Conversely, because the
        requirement grows with $N$ without bound, any \emph{fixed} logit gap
        $\Delta_0$ is eventually overwhelmed: $N > L_u e^{\Delta_0}$ suffices,
        and no amount of training removes the dependence on $N$. A model that
        has learned a gap appropriate for $10^3$ tokens has not thereby learned
        one appropriate for $10^7$.
\end{enumerate}
\end{remark}

\begin{remark}[The practical recipe]
\label{rem:recipe}
Corollary~\ref{cor:cap} inverts into a design rule. Given a target share
$1-\varepsilon$ and an estimated gap $\Delta$ for the model in use, the mutable
budget should satisfy
\begin{equation}
  C \;\le\; L_u e^{\Delta}\left(\frac{1}{\varepsilon} - 1\right).
  \label{eq:recipe}
\end{equation}
For $\varepsilon = 0.5$ this is $C \le L_u e^{\Delta}$, the half-mass threshold.
The formula is exact for the two-group model and requires only that the planner
either measures $\Delta$ or chooses a conservative lower bound for it. This is
the form in which the theory is actually usable, and \S\ref{sec:empirical}
supplies the measurement for one model.
\end{remark}

\begin{remark}[Where the model is wrong]
\label{rem:limits}
The two-group assumption is a simplification and we do not want to hide it.
Real key sets have dispersed scores within each group, and $\Delta$ as defined
above is the gap between group \emph{means}, not a gap between actual logits.
Three things follow. (i) Equation~\eqref{eq:threshold} is exact only for
constant within-group scores; otherwise it defines an effective gap whose value
is set by the observed share, which is how we estimate it below. (ii) The
worst-case theorem of \S\ref{sec:bound} is still available for the adversarial
case, and it is the reason we keep it. (iii) Multi-head structure is ignored:
$\Ainv$ is defined per head, and a claim about the layer as a whole requires an
aggregation convention, which we state when reporting measurements.
\end{remark}

% ---------------------------------------------------------------------------
\section{The Softmax Inundation Bound}
\label{sec:bound}

We include the norm-based result for completeness and for the adversarial case,
with its limitations attached.

\begin{assumption}[Bounded projection energy]
\label{assump:bounded}
There exist $M_Q, M_K > 0$ such that $\|\mathbf{q}_i\|_2 \le M_Q$ and
$\|\mathbf{k}_j\|_2 \le M_K$ for all $i,j$. By Cauchy--Schwarz,
\begin{equation}
  -B \le s_{i,j} \le B, \qquad B = \frac{M_Q M_K}{\sqrt{d_k}},
\end{equation}
and we write $\kappa = e^{2B}$.
\end{assumption}

\begin{theorem}[Softmax Inundation Bound]
\label{thm:inundation}
Under Assumption~\ref{assump:bounded}, for every query $i$ and turn $t$,
\begin{equation}
  \Ainv \;\le\; \frac{\kappa\, L_u}{L_u + N},
  \qquad\text{so}\qquad
  \Ainv = \mathcal{O}\!\left(\frac{L_u}{N}\right)
  \ \text{ as } N \to \infty .
  \label{eq:inundation}
\end{equation}
\end{theorem}

\begin{proof}
Since $e^{-B} \le e^{s_{i,j}} \le e^{B}$ for all $j$,
\begin{equation}
  \sum_{j\in\Sinv} e^{s_{i,j}} \le L_u e^{B},
\end{equation}
and dropping the non-negative $\Swork$ terms from the denominator of
\eqref{eq:attn},
\begin{equation}
  \sum_{m=1}^{N_t} e^{s_{i,m}} \ge (L_u + N)e^{-B}.
\end{equation}
Dividing gives the claim.
\end{proof}

\begin{remark}[This result is weak, and we say so]
\label{rem:weak}
Equation~\eqref{eq:inundation} is vacuous whenever
$\kappa L_u \gtrsim L_u + N$, since the right-hand side then exceeds $1$ and the
inequality carries no information. With $B = 10$, $\kappa \approx 5\times 10^{8}$;
the bound becomes informative only when $N$ is within a factor of $\kappa$ of
$L_u$, which at practical scales it is not. The result is also loose in a second
way: it bounds the share by assuming the \emph{worst} scores for the invariant
set and the \emph{best} for the distractors, whereas the observed configuration
is almost always the opposite. Its only genuine use is adversarial: it says that
without any assumption on the score ordering, volume alone cannot promise a
floor. Proposition~\ref{prop:identity} is the version to use when a gap estimate
is available.
\end{remark}

% ---------------------------------------------------------------------------
\section{Empirical measurement}
\label{sec:empirical}

\subsection{Method}

We measure attention directly rather than modelling it. The probe
(\texttt{experiments/live\_attention\_probe.py}) does the following:

\begin{enumerate}[label=(\arabic*), leftmargin=*]
  \item Loads \texttt{Qwen2.5-0.5B} in \texttt{float16} with
        \texttt{attn\_implementation="eager"}, so that attention weights are
        materialized rather than fused away.
  \item Constructs a prompt with a fixed invariant prefix (62 tokens: a system
        role, a standing objective, and four hard constraints) followed by a
        variable number of distractor tokens drawn from a set of realistic
        log-line templates, followed by a fixed task question.
  \item Runs a single forward pass with \texttt{output\_attentions=True} and
        reads the final query position's attention row over all keys.
  \item Reports the mean over all 24 layers and 14 query heads of the mass
        assigned to the invariant prefix.
\end{enumerate}

Two methodological points deserve emphasis. First, we report the
\emph{final} query position because that is the position whose distribution
governs the next generated token; averages over all positions would mix in
early positions whose visible key set is shorter. Second, we validate the
extraction against the stock implementation rather than trusting it: an earlier
version of this probe computed the final attention row from hidden states via a
hook, and we discarded that approach after it disagreed with eager attention by
$1.4\times10^{-2}$ in the maximum absolute deviation. The reported numbers come
from the unmodified library path.

\begin{remark}[What ``invariant share'' is not]
It is not a measure of task success, and it is not a probability that the model
will follow the instruction. It is the fraction of one head's attention mass at
one position that lands on a fixed set of tokens. The reason to care is that a
share near zero makes instruction-following close to impossible in that head,
whereas a healthy share merely makes it possible.
\end{remark}

\subsection{Result 1: dilution as distractors grow}

\begin{table}[ht]
\centering
\small
\begin{tabular}{rrrrr}
\toprule
\textbf{Distractors $N$} & \textbf{Total ctx.} & \textbf{Inv.\ share} & \textbf{Uniform} & \textbf{Enrichment} \\
\midrule
0    & 81   & 0.527 & 0.765 & 0.69$\times$ \\
64   & 146  & 0.427 & 0.425 & 1.00$\times$ \\
128  & 210  & 0.422 & 0.295 & 1.43$\times$ \\
256  & 338  & 0.405 & 0.183 & 2.21$\times$ \\
512  & 593  & 0.398 & 0.105 & 3.80$\times$ \\
1024 & 1106 & 0.381 & 0.056 & 6.80$\times$ \\
2048 & 2129 & 0.388 & 0.029 & 13.33$\times$ \\
\bottomrule
\end{tabular}
\caption{Qwen2.5-0.5B, final-query attention to a fixed 62-token invariant
prefix, averaged over 24 layers and 14 heads. ``Uniform'' is the share that
would obtain if attention were spread evenly, $L_u/N_t$. ``Enrichment'' is the
ratio of the two. Distractor text is drawn from log-line templates; the task
question is fixed.}
\label{tab:sweep}
\end{table}

Two features stand out, and only one of them was expected.

\emph{Expected:} the invariant share declines as distractors accumulate, from
$0.527$ at $N=0$ to $0.388$ at $N=2048$. With $L_u$ fixed and $N$ growing, the
denominator of \eqref{eq:attn} grows and a fixed set of keys must lose share
unless its relative scores grow.

\emph{Unexpected:} the decline is very small. A $32\times$ increase in
distractor volume moves the share by about $9\%$ relative. Simultaneously the
enrichment over uniform grows from $1.0\times$ to $13.3\times$: the model is
attending to the prefix far more than chance, and increasingly so. The naive
reading of \eqref{eq:threshold} with a constant $\Delta$ would predict the share
falling to roughly $0.01$ by $N=2048$; it does not.

\subsection{Result 2: the effective gap grows logarithmically}

We can make the discrepancy quantitative. Inverting
\eqref{eq:threshold} for the effective gap,
\begin{equation}
  \hat{\Delta}(N)
  = \ln\!\left(\frac{L_u}{N}\left(\frac{1}{\Ainv} - 1\right)\right),
  \label{eq:delta-hat}
\end{equation}
which is well defined whenever $\Ainv \in (0,1)$.

\begin{table}[ht]
\centering
\small
\begin{tabular}{rrrrr}
\toprule
\textbf{$N$} & \textbf{Inv.\ share} & \textbf{$\hat{\Delta}$ (nats)} & \textbf{$\ln(N/L_u)$} & \textbf{$L_u e^{\hat{\Delta}}$} \\
\midrule
64   & 0.427 & $-0.264$ & 0.032 & 48 \\
128  & 0.422 & $0.410$  & 0.725 & 93 \\
256  & 0.405 & $1.033$  & 1.418 & 174 \\
512  & 0.398 & $1.696$  & 2.111 & 338 \\
1024 & 0.381 & $2.321$  & 2.804 & 631 \\
2048 & 0.388 & $3.042$  & 3.497 & 1299 \\
\bottomrule
\end{tabular}
\caption{Effective logit gap recovered from the measured share via
\eqref{eq:delta-hat}. The final column is the half-mass threshold
$L_u e^{\Delta}$ from Corollary~\ref{cor:half}: below this $N$ the invariant set
holds the majority of attention.}
\label{tab:gap}
\end{table}

Regressing $\hat{\Delta}$ on $\ln N$ gives
\begin{equation}
  \hat{\Delta}(N) \;\approx\; 0.945\,\ln N - 4.195,
  \qquad R^2 = 0.9996 .
  \label{eq:fit}
\end{equation}

The slope is $0.945$, and the theoretical boundary value is $1$. A slope of
exactly $1$ would mean
\begin{equation}
  \hat{\Delta}(N) = \ln N + c
  \quad\Longrightarrow\quad
  \frac{N}{L_u}e^{-\hat{\Delta}} = \frac{e^{-c}}{L_u} = \text{const},
\end{equation}
so that by \eqref{eq:threshold} the invariant share is \emph{exactly
constant} in $N$. The measured slope sits just below $1$, which is why the
observed share drifts slowly downward rather than being flat: the model
compensates for context growth at $94.5\%$ of the rate required to hold its
allocation, and the residual $5.5\%$ is what shows up as dilution.

\begin{observation}
\label{obs:compensation}
On Qwen2.5-0.5B, over $N \in [64, 2048]$ and for this distractor distribution,
the invariant attention share is an approximately scale-invariant function of
distractor volume. The model supplies a logit gap that grows essentially as
$\ln N$, matching the requirement identified in Remark~\ref{rem:log} to within
a small constant margin.
\end{observation}

This is not, in our reading, a coincidence of a small model. Any training
objective that rewards instruction-following across a range of context lengths
will select for representations whose invariant-versus-background gap grows with
length, because Remark~\ref{rem:log} shows that such growth is available at
logarithmic cost. Instruction tuning across varied sequence lengths is, in
effect, a curriculum for exactly this property.

\subsection{Result 3: capping stabilises what drift remains}

If the model compensates almost perfectly, the residual drift is small, and the
question becomes whether it is worth controlling. We test the turn-scale
behaviour by simulating an append-only history against a capped one.

\begin{table}[ht]
\centering
\small
\begin{tabular}{llrr}
\toprule
\textbf{Turn} & \textbf{Arm} & \textbf{Total ctx.} & \textbf{Inv.\ share} \\
\midrule
2  & append-only   & 202  & 0.423 \\
2  & bounded ($C{=}240$) & 202  & 0.423 \\
5  & append-only   & 381  & 0.405 \\
5  & bounded       & 322  & 0.406 \\
10 & append-only   & 682  & 0.394 \\
10 & bounded       & 322  & 0.406 \\
20 & append-only   & 1282 & 0.388 \\
20 & bounded       & 322  & 0.406 \\
\bottomrule
\end{tabular}
\caption{Turn-scaled trace, 60 distractor tokens per turn. The bounded arm holds
mutable content at the cap, so total context plateaus and the share is constant
thereafter. The append-only arm drifts down by about $8\%$ relative over 20
turns.}
\label{tab:turns}
\end{table}

Bounding the mutable tiers arrests the residual drift and, more importantly,
makes it \emph{provably} bounded rather than empirically small. That is the case
for the guarantee in \S\ref{sec:guarantee}, and after Result~2 it is the honest
case: not ``we prevent collapse'', but ``we replace a property we observed with a
property the controller enforces''.

\subsection{Threats to validity}

We list these before the reader has to.

\begin{enumerate}[label=(\alph*), leftmargin=*]
  \item \textbf{One model, one size.} All measurements are Qwen2.5-0.5B. Whether
        the $0.945$ slope holds at 7B, 70B, or in models with different
        positional encodings is untested. The qualitative prediction from
        Remark~\ref{rem:log} --- slope at most~$1$, and higher for models
        trained at longer lengths --- is testable and not tested.
  \item \textbf{Synthetic distractors.} The distractor text is drawn from eight
        log-line templates with a varying integer. Real tool output is far more
        heterogeneous. Homogeneity may make the group-mean gap
        (and hence $\hat{\Delta}$) unusually well defined; a distribution with
        a few very high-scoring distractors could behave differently, and
        Remark~\ref{rem:limits} of \S\ref{sec:bound} names that case as
        adversarial and out of scope for a volume argument.
  \item \textbf{Attention is a proxy.} We measure allocation, not behaviour. The
        claim ``allocation is preserved'' does not entail ``the instruction is
        followed''. Correlating share with task success requires a behavioural
        benchmark, which we have not run.
  \item \textbf{Final-position attention only.} We read one row per head. This
        is the position that governs generation, but a model may route the
        relevant information through earlier positions or through the residual
        stream in ways a single row does not reveal.
  \item \textbf{GPU-limited reach.} The turn-scale experiment runs at 60
        tokens per turn. The original configuration of 120 tokens per turn
        exceeded the 4\,GB of available device memory in eager mode at 40 turns.
        Absolute numbers at realistic per-turn volumes will differ, though the
        qualitative plateau-versus-drift contrast should not.
  \item \textbf{Absence of collapse may be specific to the setting.} We do not
        observe inundation collapse, and we have not established the conditions
        under which it would occur. It is possible that a longer horizon than we
        can test on this hardware, or a distractor distribution with higher score
        variance, would produce it.
\end{enumerate}

\subsection{What the measurement actually settles}

\begin{table}[ht]
\centering
\small
\begin{tabular}{p{0.44\textwidth}p{0.5\textwidth}}
\toprule
\textbf{Claim} & \textbf{Status after measurement} \\
\midrule
Volume drives allocation through the softmax denominator &
\textbf{Supported.} Share is monotone decreasing in $N$ across the tested range, in the direction \eqref{eq:threshold} requires. \\
\addlinespace
A fixed logit gap is eventually overwhelmed by $N$ &
\textbf{Logically settled, not observed.} Follows from \eqref{eq:threshold}; the measured $N$ range is too small to exhibit it. \\
\addlinespace
Volume control holds the share fixed &
\textbf{Supported.} Table~\ref{tab:turns}: the capped arm is constant to three decimals. \\
\addlinespace
Unbounded volume causes attention collapse in practice &
\textbf{Not supported.} Share fell $\approx 9\%$ over a $32\times$ growth in $N$. The naive model over-predicts dilution substantially. \\
\addlinespace
Volume control improves task success &
\textbf{Untested.} Requires a behavioural benchmark. \\
\bottomrule
\end{tabular}
\caption{Status of the paper's claims after \S\ref{sec:empirical}. We include
the negative entries because they are the ones a reader is entitled to know.}
\label{tab:claims}
\end{table}

% ---------------------------------------------------------------------------
\section{Generational Context}
\label{sec:gc}

\subsection{Design}

If \S\ref{sec:threshold} identifies volume as the controllable variable, the
response is to control it explicitly. Generational Context applies the shape of
generational garbage collection and of virtual-memory paging at the granularity
of semantic context entries rather than KV blocks.

\begin{definition}[Generational Context]
\label{def:gc}
Active context is partitioned into three tiers:
\begin{enumerate}[label=\textbf{Gen-\arabic*}, leftmargin=*]
  \item \emph{Tenured.} Founding objectives, system invariants, and hard
        constraints. Pinned; never evicted. Occupies $L_u$ tokens.
  \item \emph{Evaporative.} Working memory: hypotheses, plans, uncommitted
        deltas, and committed-but-revisable state. Paged under a strict capacity
        $C_1$ by Attention-Weighted Eviction (\S\ref{sec:awe}).
  \item \emph{Ephemeral.} Raw tool output and logs. Passed through a
        distillation map $\sigma_{\mathrm{distill}}$ producing at most $K_0$
        tokens, with the raw form discarded under a turn-based TTL.
\end{enumerate}
Let $C_{\mathrm{overhead}}$ bound the structural tokens (role tags, delimiters,
section headers) needed to render the window, and set
\begin{equation}
  C_{\mathrm{bound}} = C_1 + K_0 + C_{\mathrm{overhead}} .
  \label{eq:cbound}
\end{equation}
\end{definition}

\subsection{Why three tiers and not one budget}

A single global budget would satisfy the volume bound of
Theorem~\ref{thm:guarantee} equally well. The tiering earns its place on a
different axis: it makes the eviction \emph{policy} well posed. Under a single
budget, a diagnostic observation and a standing constraint compete on the same
footing, and attention-based scoring will sometimes evict the constraint,
because in a given turn the constraint may genuinely be less attended than a
salient error. Pinning removes that possibility for the tokens where it is
intolerable, and the volume cap then bounds everything else.

The design is therefore: one tier that is exempt from competition, one tier that
competes under a budget, and one tier that is summarized before it ever
competes.

\subsection{Interaction with the threshold result}

Combining Definition~\ref{def:gc} with Remark~\ref{rem:log} gives a concrete
sizing rule. With $C_{\mathrm{bound}}$ as the mutable budget,
Corollary~\ref{cor:cap} becomes
\begin{equation}
  \Ainv \;\ge\; \frac{L_u e^{\Delta}}{L_u e^{\Delta} + C_{\mathrm{bound}}},
\end{equation}
so a designer targeting $\Ainv \ge 1-\varepsilon$ chooses
$C_{\mathrm{bound}} \le L_u e^{\Delta}(1/\varepsilon - 1)$. Nothing in this
depends on the horizon, and nothing in it depends on the model continuing to
compensate for context growth: the bound is maintained by the controller, not by
the model's properties.

% ---------------------------------------------------------------------------
\section{Attention-Weighted Eviction}
\label{sec:awe}

\subsection{Selection policy}

The guarantee of \S\ref{sec:guarantee} presumes a controller that maintains the
volume ceiling. We implement eviction as an adaptation of ARC
\cite{megiddo2003arc}, which already solves recency-versus-frequency allocation
for paged memory.

The cache maintains a recency list $T_1$, a frequency/saliency list $T_2$, ghost
registries $B_1, B_2$ recording recently evicted entries, and an adaptation
parameter $p$ targeting the size of $T_1$. Victim selection within a list uses
the \emph{Information Density Score}
\begin{equation}
  \mathrm{IDS}(E)
  = \frac{\mathrm{AttnWeight}(E)\cdot 100}{\mathrm{Entropy}(E) + 10^{-4}},
  \label{eq:ids}
\end{equation}
with tier-based priors for entries whose attention statistics have not yet been
measured. The intuition is that an entry is worth keeping when it receives
attention predictably rather than diffusely: a high weight spread across a
dispersed distribution is closer to noise than the same weight concentrated on a
few tokens.

\subsection{Post-condition}

The invariant the controller enforces is
\begin{equation}
  \sum_{e \in T_1 \cup T_2} \mathrm{tokens}(e) \;\le\; C
  \qquad \text{after every operation},
  \label{eq:post}
\end{equation}
with any single entry larger than $C$ truncated before insertion. Equation
\eqref{eq:post} is a post-condition rather than a target: it holds after
insertion, after promotion, and after distillation.

\subsection{Two correctness traps}

Both of these were found by testing rather than by inspection, and both are the
kind of error that a bound like \eqref{eq:post} would hide if it were assumed
rather than checked.

\begin{enumerate}[label=(\alph*), leftmargin=*]
  \item \textbf{Self-eviction on promotion.} Promoting an entry from $T_1$ to
        $T_2$ can grow the footprint, which triggers eviction, and a naive
        victim-selection rule can choose the newly promoted entry as its own
        victim. The entry is then silently lost in the same operation that
        promoted it, and because the operation reports success the loss is
        invisible. The fix is to exclude the incoming key from the victim pool:
        the routine must claim space from \emph{other} entries or fail.
  \item \textbf{Unbounded ghost registries.} The ghost lists record recently
        evicted keys so that re-admission can be detected. If they are not
        themselves bounded, the bookkeeping grows without limit even though the
        active cache does not, which defeats the purpose of the controller while
        leaving \eqref{eq:post} satisfied.
\end{enumerate}

% ---------------------------------------------------------------------------
\section{The Generational Retention Guarantee}
\label{sec:guarantee}

\begin{theorem}[General retention guarantee]
\label{thm:guarantee}
Fix a query $i$ and let $\Delta$ be the effective logit gap between the
invariant set and the non-invariant set. If the active non-invariant volume is
maintained at most $C_{\mathrm{bound}}$ at every turn, then for all $t \ge 1$,
\begin{equation}
  \Ainv^{(t)}
  \;\ge\;
  \frac{L_u e^{\Delta}}{L_u e^{\Delta} + C_{\mathrm{bound}}}
  \;>\; 0 .
  \label{eq:guarantee}
\end{equation}
In particular the bound is independent of the horizon $t$.
\end{theorem}

\begin{proof}
By Definition~\ref{def:gc} the non-invariant count satisfies
$N_t \le C_1 + K_0 + C_{\mathrm{overhead}} = C_{\mathrm{bound}}$ at every turn,
independently of $t$. Apply Corollary~\ref{cor:cap} with $C = C_{\mathrm{bound}}$.
The right-hand side is positive because $L_u e^{\Delta} > 0$ and
$C_{\mathrm{bound}} < \infty$.
\end{proof}

\begin{corollary}[No horizon decay]
\label{cor:no-decay}
Under the hypotheses of Theorem~\ref{thm:guarantee},
$\inf_{t\ge1} \Ainv^{(t)} > 0$.
\end{corollary}

\begin{remark}[Worst-case variant]
The same conclusion holds without any assumption on the gap, at the cost of
sharpness. Replacing $\Delta$ by the norm bound of
Theorem~\ref{thm:inundation} gives
\begin{equation}
  \Ainv^{(t)} \;\ge\; \frac{L_u}{L_u + \kappa\, C_{\mathrm{bound}}},
\end{equation}
which is positive, bounded away from zero, and horizon-independent, but
numerically uninformative at realistic $\kappa$. We state it only to record that
the qualitative conclusion --- no horizon decay --- does not depend on the gap
being favourable. The controller prevents \emph{collapse}; the gap determines
how much share is actually retained.
\end{remark}

\begin{remark}[What the guarantee is not]
It is a statement about attention allocation under a maintained volume bound. It
says that if the controller holds, then volume alone cannot drive the invariant
share to zero. It does not say the share is large, that the model will follow the
instruction, or that the task will succeed. An adversary who controls distractor
\emph{content} rather than count can violate the spirit of the result while
leaving the bound intact, by constructing distractors whose scores dominate the
invariant set within any fixed budget.
\end{remark}

% ---------------------------------------------------------------------------
\section{Intent auditing}
\label{sec:cic}

The architecture includes a component for checking that a proposed action is
consistent with the declared objective, and for leaving a verifiable trail of
those checks. It is an engineering component, and we describe it as such.

An intent is serialized as a canonical digest over the declared goal, values,
constraints, and priorities, and signed with Ed25519. A proposal is accepted only
if its authorization token verifies against the current intent digest and its
signatory holds the required role. Decisions are recorded in an append-only
hash-chained log, so that a later dispute can be resolved by recomputation
rather than by trusting the agent's account of what it decided.

Two honest caveats. The threat model is undocumented: we have not specified what
the adversary is assumed to control, and without that the component should be
read as a consistency mechanism rather than a security boundary. And the
scenarios that exercise it are authored by the same team that wrote the checker,
which is exactly the configuration in which a checker appears more capable than
it is. Both gaps are addressable, and neither is addressed here.

% ---------------------------------------------------------------------------
\section{Reference implementation}
\label{sec:impl}

The architecture is implemented in the \texttt{generational\_context} package
and depends only on the Python standard library at runtime. It exposes tool
observation, reasoning recording, and prompt compilation; it enforces the
capacity post-condition \eqref{eq:post} and the Gen-0 TTL behaviour.

The repository also contains:

\begin{itemize}[leftmargin=*]
  \item \texttt{ccb\_bench.py}, the deterministic simulation harness reported in
        the companion paper \cite{sanyal2026softmax};
  \item \texttt{experiments/live\_attention\_probe.py}, the measurement script
        behind \S\ref{sec:empirical}, with its outputs committed as CSV;
  \item \texttt{tests/}, which exercises the tier invariants, the eviction
        post-condition including the self-eviction case, and the intent-audit
        primitives. The cryptographic tests skip rather than fail when the
        optional \texttt{cryptography} dependency is absent, matching the
        library's documented fallback.
\end{itemize}

None of this is a substitute for the behavioural experiment described in
\S\ref{sec:predictions}, and it is not offered as one.

% ---------------------------------------------------------------------------
\section{Falsifiable predictions}
\label{sec:predictions}

A theory that cannot be wrong is not useful. The following are stated so they
can be tested and, if the theory is mistaken, refuted.

\begin{enumerate}[label=\textbf{P\arabic*.}, leftmargin=*]
  \item \textbf{Slope bound.} For any instruction-tuned model, the regression
        slope of $\hat{\Delta}$ on $\ln N$ is at most $1$ in the regime where
        the invariant prefix is a plausible instruction block. A slope
        materially above $1$ would mean the model \emph{over}-compensates as
        context grows, which our account of the mechanism does not predict.
        \emph{Test:} repeat \S\ref{sec:empirical} across model families and
        sizes.

  \item \textbf{Slope scales with training length.} Models trained or
        instruction-tuned at longer context lengths should exhibit slope closer
        to $1$ than models tuned at shorter lengths, because longer training
        makes the logarithmic compensation more available.
        \emph{Test:} regress slope against stated training context length across
        a model family.

  \item \textbf{Distribution sensitivity.} Replacing homogeneous log-line
        distractors with heterogeneous real tool traces should lower
        $\hat{\Delta}$ at fixed $N$, because a heterogeneous set produces higher
        maximum scores against the invariant block.
        \emph{Test:} rerun the sweep with recorded agent traces in place of
        templates.

  \item \textbf{Cap constant.} In a capped arm the measured share should be
        constant in the number of turns to within measurement noise, and should
        track \eqref{eq:cap} as $C$ varies.
        \emph{Test:} sweep $C$ and compare against the closed form.

  \item \textbf{Behavioural coupling.} Invariant share at the final position
        should correlate with instruction adherence in a multi-turn benchmark,
        weakly but positively. A null correlation would not refute the
        allocation mathematics, but it would remove the practical reason to care
        about it.
        \emph{Test:} run a long-horizon agent benchmark with the controller on
        and off, reporting both share and success rate.
\end{enumerate}

Prediction P5 is the one that decides whether this line of work matters. The
mathematics in \S\S\ref{sec:threshold}--\ref{sec:guarantee} is, we believe,
correct; the question of whether preserved allocation improves outcomes is
empirical and open, and it is the experiment we most want to see run.

% ---------------------------------------------------------------------------
\section{Limitations}
\label{sec:limits}

We collect the limitations stated throughout, in one place.

\begin{enumerate}[label=(\arabic*), leftmargin=*]
  \item \textbf{The name is not a claim.} ``Fool's Theorem'' denotes a failure
        mode. The mathematical content of this paper is
        Proposition~\ref{prop:identity}, Theorem~\ref{thm:inundation}, and
        Theorem~\ref{thm:guarantee}, none of which is novel in technique.
        Proposition~\ref{prop:identity} is a rearrangement of the softmax
        definition; Theorem~\ref{thm:inundation} is Cauchy--Schwarz.

  \item \textbf{The strongest result is conditional on a gap estimate.} The
        usable bound \eqref{eq:cap} needs $\Delta$, which is model-,
        distribution-, layer-, and head-dependent. We measured it for one model
        on one distractor distribution. Treating it as a constant across
        settings would be a misuse.

  \item \textbf{The empirical evidence partly undercuts the motivation.} As
        Result~2 shows, the model we tested already compensates for context
        growth at roughly the required rate, and the dilution we measured is
        mild. The remedy is justified by making a property enforceable rather
        than by rescuing an observed collapse. Presenting it as the latter would
        misrepresent the data.

  \item \textbf{Single head, single model, synthetic distractors, single
        position.} See \S\ref{sec:empirical}(a)--(d).

  \item \textbf{No behavioural evidence.} Nothing here shows that any of this
        improves task outcomes. It is the least interesting limitation to state
        and the most important.

  \item \textbf{Adversarial content is out of scope.} Volume control bounds
        count, not score. A bounded number of high-scoring distractors can
        dominate attention, and no result in this paper forbids it.

  \item \textbf{The intent-audit component is unvalidated} and has no
        documented threat model.
\end{enumerate}

% ---------------------------------------------------------------------------
\section{Related work}
\label{sec:related}

\paragraph{Attention sinks and streaming inference.} Xiao et al.
\cite{xiao2023streamingllm} observe that a small number of initial tokens
absorb disproportionate attention, and exploit this to stream indefinitely. Our
Observation~\ref{obs:compensation} is consistent with a related effect: the
prefix's advantage grows with context length rather than remaining fixed, which
is what permits near-scale-invariant allocation. The two accounts differ in aim
--- theirs addresses degenerate KV statistics for continuity of generation,
ours asks whether declared invariants keep competing --- but they are likely
describing overlapping machinery, and we would not be surprised if the sink
mechanism is the substrate of the compensation we measure.

\paragraph{Position sensitivity.} \emph{Lost in the Middle}
\cite{liu2023lost} documents that content in the interior of a long context is
used less effectively than content at the edges; RULER \cite{hsieh2024ruler}
shows nominal window size overstates usable context. Both are compatible with
our account and neither is subsumed by it: position and volume are distinct
variables, and our measurement holds position fixed (the invariant prefix is
always first) while varying volume.

\paragraph{KV-cache eviction.} H2O \cite{zhang2023h2o} and SnapKV
\cite{li2024snapkv} select which KV entries to retain under a memory budget,
using attention statistics as the salience signal. Attention-Weighted Eviction
uses the same signal, and IDS \eqref{eq:ids} is a close relative of their
scoring. The difference is granularity: they operate on KV entries within a
forward pass to reduce memory, whereas we operate on semantic context entries
across turns to bound a prompt. A reader who knows that literature should regard
our eviction component as an application, not a contribution.

\paragraph{Agent memory.} MemGPT \cite{packer2023memgpt} manages a virtual
context window through function calls, with paging decisions made by the model
itself. Generational Context makes those decisions with a deterministic
controller, trading flexibility for predictable cost and for a capacity
invariant that can be checked rather than hoped for. Pinning invariants
corresponds to MemGPT's always-in-context core memory; the differences are the
explicit tiering and the TTL treatment of raw observations.

\paragraph{Cache replacement.} ARC \cite{megiddo2003arc} supplies the adaptation
mechanism for recency versus frequency. We apply it unchanged in structure.

\paragraph{Attention.} The underlying mechanism is the scaled dot-product
attention of Vaswani et al. \cite{vaswani2017attention}; the rotary positional
encoding used by the measured model is due to Su et al. \cite{su2024roformer}.

% ---------------------------------------------------------------------------
\section{Conclusion}
\label{sec:conclusion}

We set out to explain a failure mode and ended up with a sharper account of its
mechanism than we started with, and a weaker claim about its severity.

The mechanism is the softmax denominator, and it is best expressed as an
identity rather than a bound: the invariant attention share equals
$\sigma(\Delta + \ln(L_u/N))$, so the model must supply a logit gap that grows as
$\ln N$ merely to hold its allocation constant. That logarithm is the whole
phenomenon. It explains why compensation is cheap and therefore likely, and it
explains why no fixed gap survives unbounded volume.

Measuring this on a real transformer produced the result we did not expect. The
model's effective gap grows at a slope of $0.945$ against $\ln N$, with
$R^2 = 0.9996$ --- almost exactly the rate required for scale-invariant
allocation. Dilution is therefore far milder in practice than the uniform-score
analysis suggests: about $9\%$ over a $32\times$ increase in distractor volume.
We report this because it is the strongest evidence against our own framing, and
because a reader who discovers it independently would be right to ask why we
omitted it.

The systems contribution survives this, but it changes character. Generational
Context does not rescue a collapsing share; it replaces a property we happened
to observe with one the controller enforces. The retention guarantee is
horizon-independent, it is stated in terms of a measurable gap, and it comes
with an implementation whose capacity invariant is checked rather than assumed.
Whether that enforced property buys anything in task outcomes is the open
question, and \S\ref{sec:predictions} states what would settle it.

% ---------------------------------------------------------------------------
\section*{Acknowledgments}

The term \emph{Artificial Civilization Intelligence} (ACI) was coined by Raja
Dharma Tej Maddala. The architecture, proofs, implementation, measurements, and
simulation harness in this document were produced by Team Auralis. The
measurement in \S\ref{sec:empirical} was run on a single consumer GPU; we thank
the reader in advance for running it on better hardware.

% ---------------------------------------------------------------------------
\begin{thebibliography}{9}

\bibitem{hsieh2024ruler}
C.-P. Hsieh, S. Sun, S. Kriman, S. Acharya, D. Rekesh, F. Jia, Y. Zhang, and
B. Ginsburg.
\newblock RULER: What's the real context size of your long-context language
  models?
\newblock \emph{arXiv:2404.06654}, 2024.

\bibitem{li2024snapkv}
Y. Li, Y. Huang, B. Yang, B. Venkitesh, A. Locatelli, H. Ye, T. Cai, P. Lewis,
  and D. Chen.
\newblock SnapKV: LLM knows what you are looking for before generation.
\newblock \emph{arXiv:2404.14469}, 2024.

\bibitem{liu2023lost}
N.~F. Liu, K. Lin, J. Hewitt, A. Paranjape, M. Bevilacqua, F. Petroni, and
P. Liang.
\newblock Lost in the middle: How language models use long contexts.
\newblock \emph{Transactions of the ACL}, 2024. arXiv:2307.03172.

\bibitem{megiddo2003arc}
N. Megiddo and D.~S. Modha.
\newblock ARC: A self-tuning, low overhead replacement cache.
\newblock In \emph{Proceedings of FAST}, 2003.

\bibitem{packer2023memgpt}
C. Packer, S. Wooders, K. Lin, V. Fang, S.~G. Patil, I. Stoica, and J.~E.
  Gonzalez.
\newblock MemGPT: Towards LLMs as operating systems.
\newblock \emph{arXiv:2310.08560}, 2023.

\bibitem{sanyal2026softmax}
S. Sanyal.
\newblock The Softmax Inundation Bound: Attention Dilution in Long-Context
  Transformers and Generational Context as a Volume-Bounding Remedy.
\newblock \emph{Team Auralis technical report}, 2026.
  \texttt{docs/papers/The\_Softmax\_Inundation\_Theorem.pdf}.

\bibitem{su2024roformer}
J. Su, M. Ahmed, Y. Lu, S. Pan, W. Bo, and Y. Liu.
\newblock RoFormer: Enhanced transformer with rotary position embedding.
\newblock \emph{Neurocomputing}, 2024. arXiv:2104.09864.

\bibitem{vaswani2017attention}
A. Vaswani, N. Shazeer, N. Parmar, J. Uszkoreit, L. Jones, A.~N. Gomez,
  L. Kaiser, and I. Polosukhin.
\newblock Attention is all you need.
\newblock In \emph{Advances in Neural Information Processing Systems}, 2017.

\bibitem{xiao2023streamingllm}
G. Xiao, Y. Tian, B. Chen, S. Han, and M. Lewis.
\newblock Efficient streaming language models with attention sinks.
\newblock In \emph{Proceedings of ICLR}, 2024. arXiv:2309.17453.

\bibitem{zhang2023h2o}
Z. Zhang, Y. Sheng, T. Zhou, T. Chen, L. Zheng, R. Cai, Z. Song, Y. Tian,
  C. R\'e, C. Barrett, Z. Wang, and B. Li.
\newblock H2O: Heavy-hitter oracle for efficient generative inference of large
  language models.
\newblock In \emph{Advances in Neural Information Processing Systems}, 2023.
  arXiv:2306.14048.

\end{thebibliography}

\end{document}
